\documentclass{article}

\usepackage{amsmath,amssymb}
\usepackage{xurl}
\usepackage[hidelinks]{hyperref}
\setlength{\emergencystretch}{2em}

\input{command}

\title{Wolf Theorem~8.6 --- the Doubly-Substochastic Half of Birkhoff}
\author{}
\date{2026-08-09}

\begin{document}
\maketitle
\gapnote{scope-restriction}{open}

\section{The source statement}

A nonnegative $d\times d$ matrix is doubly stochastic when every row sum and
every column sum equals $1$. It is doubly substochastic when every row sum and
every column sum is at most $1$. A permutation matrix has exactly one entry
equal to $1$ in each row and column. A partial permutation matrix is a
$\{0,1\}$-matrix with at most one entry equal to $1$ in each row and column.

Wolf's Theorem~8.6, called ``Birkhoff's theorem'', states two
characterizations in one theorem environment
\cite[Chapter~8, Theorem~8.6]{Wolf2012Quantum}. The local source is
\path{Notes/WolfNoteTexSource/ch08_distance_measures.tex},
lines~263--266.
\begin{enumerate}
    \item[\textbf{DS}] The set of $d\times d$ doubly stochastic matrices is
        the convex hull of the $d\times d$ permutation matrices, and its
        extreme points are exactly the permutation matrices.

    \item[\textbf{Sub}] The set of $d\times d$ doubly substochastic matrices
        is the convex hull of the $d\times d$ partial permutation matrices,
        and its extreme points are exactly the partial permutation matrices.
\end{enumerate}

\section{Formalized component}

The doubly-stochastic part \textbf{DS} is formalized by two declarations:
\begin{itemize}
    \item \leanid{TNLean.wolf_birkhoff_doublyStochastic_convexHull} states
        that the doubly stochastic matrices form the convex hull of the
        permutation matrices;

    \item \leanid{TNLean.wolf_birkhoff_extremePoints_doublyStochastic} states
        that the extreme points of the doubly stochastic matrices are
        exactly the permutation matrices.
\end{itemize}
Both declarations are in
\doclink{QICLean/Analysis/Birkhoff}. They directly reuse Mathlib's
\leanid{doublyStochastic_eq_convexHull_permMatrix} and
\leanid{extremePoints_doublyStochastic}.

\section{Missing component}

The doubly-substochastic part \textbf{Sub} of Theorem~8.6
\cite{Wolf2012Quantum} is not formalized. A complete treatment requires the
following definitions and results.
\begin{enumerate}
    \item Define the predicate ``doubly substochastic'' for a nonnegative
        square matrix whose row and column sums are bounded by $1$.

    \item Define partial permutation matrices as square
        $\{0,1\}$-matrices with at most one $1$ in each row and each column.

    \item Prove that the doubly substochastic matrices form the convex hull
        of the partial permutation matrices and that their extreme points are
        exactly the partial permutation matrices.
\end{enumerate}

The substochastic polytope contains the stochastic polytope as a face, but its
support may omit rows or columns. Its extreme-point proof must therefore
replace permutation matrices by partial permutation matrices.

\section{Elimination plan}

First, adapt Mathlib's \leanid{doublyStochastic} predicate to a predicate
\leanid{doublySubstochastic}, replacing the row-sum and column-sum equalities
by upper bounds. Then define partial permutation matrices over
$\operatorname{Fin}d$ and prove that they are doubly substochastic.

For the convex-hull statement, let $M$ be a $d\times d$ doubly
substochastic matrix and define
\begin{align}
    u_i
        &=\sum_j M_{ij},
    &
    v_j
        &=\sum_i M_{ij},
    \label{eq:wolf_ch8_birkhoff_doubly_substochastic_gap_row_column_sums}\\
    r_i
        &=1-u_i,
    &
    c_j
        &=1-v_j,
    \label{eq:wolf_ch8_birkhoff_doubly_substochastic_gap_row_column_deficits}\\
    s
        &=\sum_{i,j}M_{ij}.
    \label{eq:wolf_ch8_birkhoff_doubly_substochastic_gap_total_mass}
\end{align}
If $s>0$, define $S$ by
\begin{align}
    S_{ji}
        &=\frac{v_ju_i}{s}.
    \label{eq:wolf_ch8_birkhoff_doubly_substochastic_gap_completion_block_positive_mass}
\end{align}
If $s=0$, set
\begin{align}
    S
        &=0.
    \label{eq:wolf_ch8_birkhoff_doubly_substochastic_gap_completion_block_zero_mass}
\end{align}
For $s>0$, the row and column sums of $S$ are
\begin{align}
    \sum_i S_{ji}
        &=v_j,
    \label{eq:wolf_ch8_birkhoff_doubly_substochastic_gap_completion_row_sum}\\
    \sum_j S_{ji}
        &=u_i.
    \label{eq:wolf_ch8_birkhoff_doubly_substochastic_gap_completion_column_sum}
\end{align}
The same identities hold when $s=0$, because then $M=0$ and all $u_i$ and
$v_j$ vanish.

Define the block matrix
\begin{align}
    \widetilde M
        &=
        \begin{pmatrix}
            M&\operatorname{diag}(r)\\
            \operatorname{diag}(c)&S
        \end{pmatrix}.
    \label{eq:wolf_ch8_birkhoff_doubly_substochastic_gap_doubly_stochastic_completion}
\end{align}
Its row and column sums satisfy
\begin{align}
    \sum_b\widetilde M_{ab}
        &=1,
    \label{eq:wolf_ch8_birkhoff_doubly_substochastic_gap_completed_row_sums}\\
    \sum_a\widetilde M_{ab}
        &=1.
    \label{eq:wolf_ch8_birkhoff_doubly_substochastic_gap_completed_column_sums}
\end{align}
Thus $\widetilde M$ is a $2d\times2d$ doubly stochastic matrix. Applying the
formalized doubly-stochastic Birkhoff theorem to $\widetilde M$ and restricting
each permutation matrix to the original $d\times d$ block yields a convex
decomposition of $M$ into partial permutation matrices.

The final step is to adapt the stochastic extreme-point argument and prove
that the extreme points of the doubly substochastic set are precisely the
partial permutation matrices.

\section{Conclusion}

This note records a scope restriction. The doubly-stochastic part of Wolf's
Theorem~8.6 \cite{Wolf2012Quantum}, including both the convex-hull and
extreme-point characterizations, is fully formalized. The definition and
polytope theory of doubly substochastic matrices, including the
partial-permutation extreme-point characterization, remain unformalized.
Tracking:
\href{https://github.com/LionSR/TNLean/issues/3959}{TNLean issue \#3959}.

\bibliographystyle{alpha}
\bibliography{references}

\end{document}
